% Generated by support/export-artifacts.py; do not edit.
% Slug: ch62-class
% Source: chapters/ch62.tex:56
% TikZ source: chapters/ch62.tex:58
% Original label: \label{fig:ch62-class}
% Original caption (LaTeX source):
%   Identity rewrites put several finite expressions in one class; recursive class references also describe longer alternatives.
% Caption: Identity rewrites put several finite expressions in one class; recursive class references also describe longer alternatives.
% Setup: support/compiler-pseudocode.sty (TikZ libraries and shared styles).
\begin{tikzpicture}[node distance=14mm]
\node[basic block] (root) {root e-class};
\node[basic block,below left=of root] (x) {x};
\node[basic block,below=of root] (add) {add(root, zero)};
\node[basic block,below right=of root] (mul) {mul(root, one)};
\draw[cfg edge] (root)--(x); \draw[cfg edge] (root)--(add); \draw[cfg edge] (root)--(mul);
\end{tikzpicture}

% Original figure LaTeX:
% \begin{figure}[htbp]
% \centering
% \begin{tikzpicture}[node distance=14mm]
% \node[basic block] (root) {root e-class};
% \node[basic block,below left=of root] (x) {x};
% \node[basic block,below=of root] (add) {add(root, zero)};
% \node[basic block,below right=of root] (mul) {mul(root, one)};
% \draw[cfg edge] (root)--(x); \draw[cfg edge] (root)--(add); \draw[cfg edge] (root)--(mul);
% \end{tikzpicture}
% \caption{Identity rewrites put several finite expressions in one class; recursive class references also describe longer alternatives.}
% \label{fig:ch62-class}
% \end{figure}
